% =====================================================================
%  ch02b.tex —— 第 2 章（后半）：原书 p.22–35（PDF 31–44）
%  2.6 相对论性点粒子 ～ 2.12 探测器的轨道
% =====================================================================

% [原文 p.22]
\begin{itemize}[leftmargin=2em]
  \item 作用量是配位空间 $W$ 上路径的函数；
  \item 存在称为 Lagrange 函数、定义在 $TW$ 上的函数 $L$；
  \item 作用量由沿路径积分 Lagrange 函数来定义，
\end{itemize}
而在哈密顿形式的分析力学中，
\begin{itemize}[leftmargin=2em]
  \item 存在称为哈密顿函数、定义在相空间 $T^*W$ 上的函数 $H$；
  \item $H$ 与 $L$ 通过 Legendre 变换相联系；
  \item 作用量是相空间 $T^*W$ 上路径的函数；
  \item 作用量由沿路径积分 1 次微分形式 $\bm{p}\cdot\rd\bm{q}-H\rd t$ 来定义。
\end{itemize}

称 1 次微分形式 $\bm{p}\cdot\rd\bm{q}$ 为 \textbf{Liouville 形式}（Liouville 1-form），称 1 次微分形式 $\bm{p}\cdot\rd\bm{q}-H\rd t$ 为 \textbf{Poincaré--Cartan 积分不变量}（Poincaré--Cartan integral invariant）。

% [Tutor Note: 背景——$\bm{p}\cdot\rd\bm{q}$ 是余切丛 $T^*W$ 上的重言 1 形式（tautological form），本译注后文亦称 Liouville 形式；其外微分给出 $T^*W$ 上的标准辛形式。]

\section{相对论性点粒子}

在 Newton 力学中，两个惯性系 $(x,t)$ 与 $(x',t')$ 之间的坐标变换由 \textbf{Galilei 变换}（Galilean transformation）
\begin{equation}
  \left\{
  \begin{aligned}
    x' &= x-vt,\\
    t' &= t
  \end{aligned}
  \right.
  \tag{2.6.1}
\end{equation}
给出；而在狭义相对论中与之对应的是 \textbf{Lorentz 变换}（Lorentz transformation）
\begin{equation}
  \left\{
  \begin{aligned}
    x' &= \frac{1}{\sqrt{1-v^2/c^2}}\,(x-vt),\\
    t' &= \frac{1}{\sqrt{1-v^2/c^2}}\,t
  \end{aligned}
  \right.
  \tag{2.6.2}
\end{equation}
作为点粒子的作用量，具有 Lorentz 不变性的泛函中最简单的是
\begin{equation}
  S = -mc\int \rd s
  \tag{2.6.3}
\end{equation}
其中
\begin{equation}
  \rd s^2 = c^2\rd t^2-\rd\bm{x}^2 = c^2\rd t^2-\rd x^2-\rd y^2-\rd z^2
  \tag{2.6.4}
\end{equation}
是微小世界间隔（infinitesimal world interval），$m$ 与 $c$ 是为调整量纲而引入的常数，其符号的取法使 $S$ 在驻点处不取最大值而取最小值。将其改写为
\begin{equation}
  S = -mc^2\int\sqrt{1-\dot{x}^2/c^2}\,\rd t
  \tag{2.6.5}
\end{equation}
% [原文 p.23]
这样一来，便得到
\begin{equation}
  L = -mc^2\sqrt{1-\dot{x}^2/c^2}
  \tag{2.6.6}
\end{equation}
正则动量为
\begin{equation}
  p = \frac{\partial L}{\partial\dot{x}} = \frac{m\dot{x}}{\sqrt{1-\dot{x}^2/c^2}}
  \tag{2.6.7}
\end{equation}
而哈密顿函数为
\begin{equation}
  H = p\cdot\dot{x}-L = \frac{mc^2}{\sqrt{1-\dot{x}^2/c^2}}
  \tag{2.6.8}
\end{equation}
此处令 $\dot{x}=0$，就得到点粒子的静止能量
\begin{equation}
  E = mc^2
  \tag{2.6.9}
\end{equation}

% [Tutor Note: 难点——(2.6.8) 中 $\bm{p}\cdot\dot{\bm{x}}=m\dot{x}^2/\sqrt{1-\dot{x}^2/c^2}$，故 $H=mc^2\bigl(1-\dot{x}^2/c^2+(\dot{x}/c)^2\bigr)/\sqrt{1-\dot{x}^2/c^2}=mc^2/\sqrt{1-\dot{x}^2/c^2}$。]

\takeaway{自由点粒子的作用量取世界间隔的线积分，其哈密顿函数给出相对论能量与动量。}

\section{Poisson 定理}

对相空间上的函数 $f(\bm{q},\bm{p})$ 与 $g(\bm{q},\bm{p})$，其 \textbf{泊松括号}（Poisson bracket）定义为
\begin{equation}
  \{f,g\} = \sum_{i=1}^{n}\left(\frac{\partial f}{\partial p_i}\frac{\partial g}{\partial q_i}-\frac{\partial g}{\partial p_i}\frac{\partial f}{\partial q_i}\right)
  \tag{2.7.1}
\end{equation}
利用它，哈密顿正则方程可改写为
\begin{equation}
  \frac{\rd\bm{p}}{\rd t} = \{H,\bm{p}\},\qquad
  \frac{\rd\bm{q}}{\rd t} = \{H,\bm{q}\}
  \tag{2.7.2}
\end{equation}
由定义容易看出，泊松括号具有下述 3 条性质：
\begin{itemize}[leftmargin=2em]
  \item 反对称性：$\{f,g\} = -\{g,f\}$．
  \item Leibniz 法则：$\{f,gh\} = \{f,g\}h+g\{f,h\}$．
  \item Jacobi 恒等式：$\{f,\{g,h\}\}+\{g,\{h,f\}\}+\{h,\{f,g\}\} = 0$．
\end{itemize}
相空间上的函数 $f(\bm{p},\bm{q})$ 随时间的发展为
\begin{align}
  \frac{\rd f}{\rd t} &= \frac{\partial f}{\partial\bm{p}}\cdot\frac{\rd\bm{p}}{\rd t}
    +\frac{\partial f}{\partial\bm{q}}\cdot\frac{\rd\bm{q}}{\rd t}
  \tag{2.7.3}\\
  &= -\frac{\partial f}{\partial\bm{p}}\cdot\frac{\partial H}{\partial\bm{q}}
    +\frac{\partial f}{\partial\bm{q}}\cdot\frac{\partial H}{\partial\bm{p}}
  \tag{2.7.4}\\
  &= \{H,f\}
  \tag{2.7.5}
\end{align}
随时间发展而不变的相空间上的函数，称为该系统的\textbf{守恒量}（conserved quantity）。这与说它与哈密顿函数的泊松括号为零是同一件事。由泊松括号的 Leibniz 法则可知，当 $f$ 与 $g$ 都是系统的守恒量时，其乘积 $fg$ 也仍是系统的守恒量。这称为 \textbf{Poisson 定理}。

% [Tutor Note: 难点——(2.7.4) 由 (2.7.2) 代入 (2.7.3) 而得：第一项 $-\partial f/\partial p_i\cdot\partial H/\partial q_i$，第二项 $+\partial f/\partial q_i\cdot\partial H/\partial p_i$，两项之和恰为 (2.7.1) 中的 $\{H,f\}$。]

\takeaway{泊松括号把正则方程统一为 $\dot{f}=\{H,f\}$，守恒量正是与 $H$ 的泊松括号为零的函数。}

% [原文 p.24]
\section{Poisson 流形}

把相空间 $T^*W$ 上的泊松括号 (2.7.1) 所具有的性质抽象化，便得到下述定义。

\begin{definition}{2.8.1}
\textbf{泊松流形}（Poisson manifold）是指流形 $M$ 与映射 $\{-,-\}\colon C^\infty(M)\times C^\infty(M)\to C^\infty(M)$ 的组 $(M,\{-,-\})$，它满足以下 3 个条件：
\begin{enumerate}[label=\arabic*.,leftmargin=2em]
  \item 反对称性：$\{f,g\} = -\{g,f\}$；
  \item Leibniz 法则：$\{f,gh\} = \{f,g\}h+g\{f,h\}$；
  \item Jacobi 恒等式：$\{f,\{g,h\}\}+\{g,\{h,f\}\}+\{h,\{f,g\}\} = 0$．
\end{enumerate}
\end{definition}

给定了泊松流形 $(M,\{-,-\})$ 及其上的函数 $H\in C^\infty(M)$，则由 $H$ 生成的 $M$ 上的 \textbf{哈密顿向量场}（Hamiltonian vector field）$X_H$ 由
\begin{equation}
  X_H = \{H,-\}
  \tag{2.8.1}
\end{equation}
确定。此向量场生成的流
\begin{equation}
  \varphi_H\colon M\times\R\to M
  \tag{2.8.2}
\end{equation}
称为以 $H$ 为哈密顿函数的 $(M,\{-,-\})$ 上的 \textbf{哈密顿流}（Hamiltonian flow）。

对泊松流形 $(M,\{-,-\})$，把 $M$ 上的函数 $f$ 送到向量场 $X_f$ 的映射
\begin{equation}
  X_{\bullet}\colon C^\infty(M)\to\Gamma(TM)
  \tag{2.8.3}
\end{equation}
是李代数的同态：

\begin{proposition*}{2.8.2}
对任意的 $f,g\in C^\infty(M)$，有
\begin{equation}
  [X_f,X_g] = X_{\{f,g\}}.
  \tag{2.8.4}
\end{equation}
\end{proposition*}

\begin{proof}
对任意的 $h\in C^\infty(M)$，
\begin{align}
  [X_f,X_g]h &= X_fX_gh-X_gX_fh
  \tag{2.8.5}\\
  &= \{f,\{g,h\}\}-\{g,\{f,h\}\}
  \tag{2.8.6}\\
  &= \{f,\{g,h\}\}+\{g,\{h,f\}\}
  \tag{2.8.7}\\
  &= -\{h,\{f,g\}\}
  \tag{2.8.8}\\
  &= \{\{f,g\},h\}
  \tag{2.8.9}\\
  &= X_{\{f,g\}}h
  \tag{2.8.10}
\end{align}
成立，故 $[X_f,X_g] = X_{\{f,g\}}$。
\end{proof}

% [Tutor Note: 难点——(2.8.6)$\to$(2.8.7) 只用内层括号的反对称性 $\{f,h\}=-\{h,f\}$；(2.8.7)$\to$(2.8.8) 才是 Jacobi 恒等式。]

\takeaway{泊松流形把泊松括号抽象为流形上的结构，哈密顿向量场 $X_f=\{f,-\}$ 生成哈密顿流。}

% [原文 p.25]
\section{Noether 定理}

在哈密顿形式的分析力学中，作为运动舞台的相空间是预先给定的；每选取其上的一个函数 $f$，就确定以它为哈密顿函数的哈密顿流 $\varphi\colon M\times\R\to M$。当哈密顿函数 $H$ 沿此流保持为一定值（即对任意 $x\in M$，$H(\varphi(x,s))$ 不依赖于 $s\in\R$）时，称 $f$ 生成的哈密顿流是以 $H$ 为哈密顿函数的哈密顿系统的\textbf{对称性}（symmetry）。$f$ 生成的哈密顿流是以 $H$ 为哈密顿函数的哈密顿系统的对称性的必要充分条件是
\begin{equation}
  \frac{\rd H}{\rd s} = \{f,H\} = 0
  \tag{2.9.1}
\end{equation}
另一方面，$f$ 是以 $H$ 为哈密顿函数的哈密顿系统的守恒量的必要充分条件是
\begin{equation}
  \frac{\rd f}{\rd t} = \{H,f\} = 0
  \tag{2.9.2}
\end{equation}
由泊松括号的反对称性，(2.9.1) 与 (2.9.2) 等价。也就是说，对称性与守恒定律是同一回事。这一事实称为 \textbf{Noether 定理}（Noether's theorem）。这个定理同时具备真正伟大的定理所应有的两个特征：
\begin{itemize}[leftmargin=2em]
  \item 主张深刻，且带来许多推论；
  \item 证明是显然的。
\end{itemize}
下面举出该定理的几个典型应用例。

% [Tutor Note: 背景——Emmy Noether（1882--1935）1918 年证明的定理把连续对称性与守恒量对应起来；此处是其在哈密顿力学中的形式，(2.9.1) 即沿 $X_f$ 的方向导数。]

\takeaway{Noether 定理断言：哈密顿流的对称性与守恒量是同一件事。}

\begin{example}{2.9.1}
以相空间上的函数为哈密顿函数的哈密顿系统称为\textbf{自励系统}（autonomous system）。系统为自励系统，与系统关于时间平移不变（即哈密顿函数不显含时间）是等价的。哈密顿函数生成的流是时间平移（即系统的时间发展），此时 Noether 定理主张：系统关于时间平移不变与能量守恒是等价的。（含时的哈密顿函数的例子将在 2.12 节看到。）
\end{example}

\begin{example}{2.9.2}
若存在某个 $i$ 使哈密顿函数 $H$ 不依赖于 $q_i$，则称 $q_i$ 为\textbf{循环坐标}。此时与 $q_i$ 正则共轭的动量 $p_i$ 成为以 $H$ 为哈密顿函数的哈密顿系统的守恒量。以 $p_i$ 为哈密顿函数的哈密顿向量场为
\begin{equation}
  X_{p_i} = \frac{\partial}{\partial q_i}
  \tag{2.9.3}
\end{equation}
注意到它生成的流是关于 $q_i$ 的平移，便知此时的 Noether 定理主张：系统关于某个坐标的空间平移不变性，与该坐标正则共轭动量的守恒定律等价。
\end{example}

% [原文 p.26]
\begin{example}{2.9.3}
考虑球坐标 $(r,\theta,\phi)$。作为例 2.9.2 的特殊情形，系统球对称，即势 $V(r,\theta,\phi)$ 只依赖于 $r$ 而不依赖于 $\theta$ 与 $\phi$，与系统的角动量 $p_\theta$ 及 $p_\phi$ 守恒是等价的。当势 $V$ 依赖于 $r$ 与 $\theta$ 而不依赖于 $\phi$ 时，在正则共轭动量中 $p_r$ 与 $p_\theta$ 不守恒，但 $p_\phi$ 守恒。
\end{example}

\section{叶状结构}

流形 $M$ 的\textbf{叶状结构}（foliation）是指把 $M$ 分割为互不相交的 $p$ 维连通子流形之族 $\mathcal{F}=\{L_\alpha\}_{\alpha\in A}$ 的分割
\begin{equation}
  M = \bigcup_{\alpha\in A}L_\alpha,\qquad
  \alpha\neq\beta\ \text{则}\ L_\alpha\cap L_\beta=\emptyset
  \tag{2.10.1}
\end{equation}
且满足：对 $M$ 的任意点 $p$，存在开邻域 $U\ni p$ 与微分同胚 $\varphi\colon U\to\R^n$，使得对任意 $\alpha\in A$ 以及 $U\cap L_\alpha$ 的任意连通分支 $F$，$\varphi(F)$ 都可由某个 $c\in\R^{n-p}$ 表为
\begin{equation}
  \varphi(F) = c\times\R^p\subset\R^n
  \tag{2.10.2}
\end{equation}
此处子流形是指单射浸入（immersion）的像。此时的 $p$ 称为叶状结构的\textbf{维数}（dimension），$\mathcal{F}$ 的元称为叶状结构的\textbf{叶}（leaf）。

叶状结构的典型例子是：对给定的无理数 $\theta\in\R\setminus\Q$，把环面 $\T^2=\R^2/\Z^2$ 分割为斜率为 $\theta$ 的直线所成的 \textbf{Kronecker 叶状结构}（Kronecker foliation）：
\[
  \T^2 = \bigcup_{[\alpha]\in\R/(\Z+\theta\Z)}L_\alpha,\qquad
  L_\alpha = \{(x,y)\in\R^2/\Z^2\mid y=\theta x+\alpha\}.
\]

% [Tutor Note: 背景——$\theta$ 无理时每条叶在 $\T^2$ 中稠密，故 Kronecker 叶状结构是「叶非闭」的典型例子，与后面的辛叶（其叶为子流形）形成对照。]

\takeaway{叶状结构把流形分割为等维数的叶，其切分布是对合的。}

流形切丛的子丛称为\textbf{分布}（distribution）。当分布 $E\subset TM$ 满足
\begin{equation}
  X,Y\in\Gamma(E)\ \text{则}\ [X,Y]\in\Gamma(E)
  \tag{2.10.3}
\end{equation}
时，称 $E$ 是\textbf{对合的}（involutive）。

对分布 $E\subset TM$，把
\begin{equation}
  \{\omega\in\Gamma(T^*M)\mid \text{对任意}\ X\in\Gamma(M,E)\ \text{有}\ \omega(X)=0\}
  \tag{2.10.4}
\end{equation}
所生成的
\begin{equation}
  A^*(M) = \bigoplus_{p=0}^{\dim M}\Gamma\Bigl(\bigwedge\nolimits^{p}T^*M\Bigr)
  \tag{2.10.5}
\end{equation}
的理想记为 $\mathcal{J}(E)$。当 $A^*(M)$ 的理想 $\mathcal{J}$ 满足 $\rd\mathcal{J}\subset\mathcal{J}$ 时，称 $\mathcal{J}$ 为\textbf{微分理想}（differential ideal）。

% [Tutor Note: 校勘——(2.10.4) 原书作 $\Gamma(M,E)$，与 (2.10.3) 的 $\Gamma(E)$ 记号不一致，疑为 $\Gamma(E)$ 之误，此处照录。]

部分流形 $N$ 若对任意 $p\in N$ 满足 $T_pN=E_p\subset T_pM$，则称 $N$ 为 $E$ 的\textbf{积分流形}（integral manifold）。在连通积分流形中关于包含关系极大的那些，称为\textbf{极大积分流形}（maximal integral manifold）。若对任意点 $p\in M$ 都存在过 $p$ 的 $E$ 的积分流形，则称 $E$ 为\textbf{可积的}（integrable）。对流形 $M$ 上可积的切分布 $E\subset TM$，$E$ 的极大积分流形之集确定了 $M$ 的叶状结构。此叶状结构的维数与 $E$ 的秩一致。

% [原文 p.27]
\begin{theorem}{2.10.1}
（Frobenius 定理（Frobenius theorem））对切分布 $E\subset TM$，下述 3 个条件等价：
\begin{enumerate}[label=\arabic*.,leftmargin=2em]
  \item $E$ 是对合的。
  \item $\mathcal{J}(E)$ 是微分理想。
  \item $E$ 是可积的。
\end{enumerate}
\end{theorem}

此定理的证明例如参见 [森田 05]。

% [Tutor Note: 背景——Frobenius 定理（1877）是叶状结构理论的基础；条件 2 中 $\rd\mathcal{J}\subset\mathcal{J}$ 即「外微分不越出理想」，与对合性互为表里。]

\section{辛流形}

对泊松流形 $(M,\{-,-\})$ 上的点 $p\in M$，由
\begin{equation}
  E_p = \{X_f(p)\in T_pM\mid f\in C^\infty(M)\}
  \tag{2.11.1}
\end{equation}
定义的各子空间构成 $TM$ 的子丛，将其记作 $E$。由命题 2.8.2，
\begin{equation}
  [X_f,X_g] = X_{\{f,g\}}
  \tag{2.11.2}
\end{equation}
故 $E$ 是对合的，从而由定理 2.10.1 知 $E$ 可积。取 $E$ 的一个极大积分流形，记作 $F$。

\begin{proposition*}{2.11.1}
对任意 $p\in F$，映射
\begin{equation}
  \gamma_p\colon T_pF\to T_p^*F,\qquad X_f\mapsto \rd(f|_F)
  \tag{2.11.3}
\end{equation}
是双射。
\end{proposition*}

\begin{proof}
由 $F$ 的定义，对任意 $g\in C^\infty(M)$ 有 $X_g\in T_pF$，故对任意 $f\in C^\infty(M)$，
% [原文 p.28]
\begin{equation}
  \rd(f|_F)(X_g) = \rd f(X_g)
  \tag{2.11.4}
\end{equation}
成立。又，$T_pF$ 的任意元都可由适当的 $g\in C^\infty(M)$ 表为 $X_g$ 的形状。由于
\begin{align}
  \rd f(X_g) &= X_gf
  \tag{2.11.5}\\
  &= \{g,f\}
  \tag{2.11.6}\\
  &= -\{f,g\}
  \tag{2.11.7}\\
  &= -X_fg
  \tag{2.11.8}
\end{align}
所以 $X_f=0$ 与 $\rd(f|_F)=0$ 等价。这说明 $\gamma_p$ 定义明确，且为单射。维数相等的向量空间之间的单射线性映射是同构，故 $\gamma_p$ 是同构。
\end{proof}

于是，对由 $M$ 的泊松结构确定的 $M$ 上叶状结构的叶，由 $M$ 的泊松结构定出的映射 $\gamma\colon TF\to T^*F$ 是向量丛的同构。利用它，对 $p\in F$ 与 $X,Y\in T_pF$，定义 $\omega\in\Gamma(T^*F\otimes T^*F)$ 为
\begin{equation}
  \omega_p(X,Y) = -\gamma_p(X)(Y)
  \tag{2.11.9}
\end{equation}
此处右端是 $-\gamma(X)\in T^*F_p$ 作用于 $Y\in T_pF$ 的像。

\begin{proposition*}{2.11.2}
$\omega$ 是反对称的；即
\begin{equation}
  \omega\in\Gamma\Bigl(\bigwedge\nolimits^{2}T^*F\Bigr).
  \tag{2.11.10}
\end{equation}
\end{proposition*}

\begin{proof}
注意对任意 $p\in F$，$T_pF$ 的任意 2 个元都可由适当的 $f,g\in C^\infty(M)$ 表为 $X_f,X_g$ 的形状，则
\begin{align}
  \omega_p(X_f,X_g) &= -\gamma(X_f)(X_g)
  \tag{2.11.11}\\
  &= -\rd f(X_g)
  \tag{2.11.12}\\
  &= -X_gf
  \tag{2.11.13}\\
  &= -\{g,f\}
  \tag{2.11.14}\\
  &= \{f,g\}
  \tag{2.11.15}\\
  &= -\omega_p(X_g,X_f)
  \tag{2.11.16}
\end{align}
故 $\omega$ 为反对称。
\end{proof}

\begin{proposition*}{2.11.3}
$\omega$ 是闭形式（closed form）。
\end{proposition*}

\begin{proof}
按外微分不依赖坐标的定义，
% [原文 p.29]
\begin{align*}
  \rd\omega(X_f,X_g,X_h)
  &= X_f\omega(X_g,X_h)-X_g\omega(X_f,X_h)+X_h\omega(X_f,X_g)\\
  &\quad -\omega([X_f,X_g],X_h)+\omega([X_f,X_h],X_g)-\omega([X_g,X_h],X_f)\\
  &= X_f(\{g,h\})-X_g(\{f,h\})+X_h(\{f,g\})\\
  &\quad -\omega([X_f,X_g],X_h)+\omega([X_f,X_h],X_g)-\omega([X_g,X_h],X_f)\\
  &= X_f(\{g,h\})-X_g(\{f,h\})+X_h(\{f,g\})\\
  &\quad -\omega(X_{\{f,g\}},X_h)+\omega(X_{\{f,h\}},X_g)-\omega(X_{\{g,h\}},X_f)\\
  &= X_f(\{g,h\})-X_g(\{f,h\})+X_h(\{f,g\})\\
  &\quad +\omega(X_h,X_{\{f,g\}})-\omega(X_g,X_{\{f,h\}})+\omega(X_f,X_{\{g,h\}})\\
  &= \{f,\{g,h\}\}-\{g,\{f,h\}\}+\{h,\{f,g\}\}\\
  &\quad +\{h,\{f,g\}\}-\{g,\{f,h\}\}+\{f,\{g,h\}\}\\
  &= 2\bigl(\{f,\{g,h\}\}+\{g,\{h,f\}\}+\{h,\{f,g\}\}\bigr)
\end{align*}
故泊松括号的 Jacobi 律与 $\omega$ 为闭形式等价。
\end{proof}

% [Tutor Note: 背景——(2.11.9) 的 $\omega$ 正是辛叶上的辛形式，此命题说明「$\rd\omega=0$」与「Jacobi 恒等式」是同一条式子的两种说法。]

将 $\omega$ 所具有的性质抽象化，得到下述定义。

\begin{definition}{2.11.4}
流形 $M$ 与其上的非退化闭 2 次微分形式 $\omega$ 的组 $(M,\omega)$ 称为\textbf{辛流形}（symplectic manifold），此时的 $\omega$ 称为\textbf{辛形式}（symplectic form）。
\end{definition}

不过，流形 $M$ 上的 2 次微分形式 $\omega\in C^\infty(M,\bigwedge^2T^*M)$ \textbf{非退化}（non-degenerate）是指：在任意点 $x\in M$，$\omega$ 所定的交错 2 次形式
\begin{equation}
  \omega_x\colon T_xM\times T_xM\to\R
  \tag{2.11.17}
\end{equation}
非退化（即：对任意 $v\in T_xM$，若 $\omega(v,-)\colon T_x\to\R$ 为零映射，则 $v=0$）。

% [Tutor Note: 校勘——(2.11.17) 之后的正文原书作 $\omega(v,-)\colon T_x\to\R$，按上下文应为 $T_xM\to\R$，此处照录。]

一般地，对向量丛的射 $\alpha\colon E\to F$，$\alpha$ 的转置给出射 $\alpha^*\colon F^*\to E^*$。给出 2 次微分形式 $\omega$，等价于给出满足反称性
\begin{equation}
  \gamma^* = -\gamma
  \tag{2.11.18}
\end{equation}
的映射
\begin{equation}
  \gamma\colon T_xM\to T_x^*M,\qquad X\mapsto\iota_X\omega
  \tag{2.11.19}
\end{equation}
而 $\omega$ 的非退化性等价于 $\gamma$ 为同构。

% [原文 p.30]
辛形式 $\omega$ 是 $\Gamma(\bigwedge^2T^*M)$ 的元；与之相仿，泊松结构 $\{-,-\}$ 可看作 $\Gamma(\bigwedge^2TM)$ 的元 $\beta$，它由
\begin{equation}
  \{f,g\} = \beta(\rd f,\rd g)
  \tag{2.11.20}
\end{equation}
刻画。此 $\beta$ 称为 \textbf{Poisson 双向量}（Poisson bivector）。非退化的泊松双向量引起 $TM$ 与 $T^*M$ 的同构，经此同构把 $\beta\in\Gamma(\bigwedge^2TM)$ 移到 $\Gamma(\bigwedge^2T^*M)$ 所得之物，就是 $M$ 的辛形式。

另一方面，辛形式按定义总是非退化的，故任意辛流形以下述方式具有泊松流形的结构：对辛流形 $(M,\omega)$ 上的函数 $f\colon M\to\R$，$f$ 生成的哈密顿向量场 $X_f$ 由
\begin{equation}
  X_f = \gamma^{-1}(\rd f)
  \tag{2.11.21}
\end{equation}
定义。由这一定义，$X_f$ 由
\begin{equation}
  \iota_{X_f}\omega = \rd f
  \tag{2.11.22}
\end{equation}
刻画。此时，任意辛流形具有由
\begin{align}
  \{f,g\} &= \omega(X_f,X_g)
  \tag{2.11.23}\\
  &= X_fg
  \tag{2.11.24}\\
  &= -X_gf
  \tag{2.11.25}
\end{align}
定出的自然泊松结构。

从而，辛流形是泊松流形的特殊情形，而泊松流形是辛流形在「非退化」这一条件上的一般化。
任意泊松流形 $M$ 都有自然的叶状结构，其每一片叶具有辛结构。这些叶称为该泊松流形的\textbf{辛叶}（symplectic leaf）。特别地，当 $M$ 是辛流形时，其作为泊松流形的辛叶就是 $M$ 全体。泊松流形上的哈密顿系统的时间发展停留在一片辛叶之上，因此在考虑一个哈密顿系统时，把舞台限定为辛流形而非泊松流形也不会失去一般性。

对辛流形 $(M_1,\omega_1)$ 与 $(M_2,\omega_2)$，若存在微分同胚 $f\colon M_1\to M_2$ 满足
\begin{equation}
  f^*\omega_2 = \omega_1
  \tag{2.11.26}
\end{equation}
则称 $(M_1,\omega_1)$ 与 $(M_2,\omega_2)$ 为辛同相（symplectomorphic），此时的 $f$ 称为\textbf{辛同相写像}（symplectomorphism）。

最基本的辛流形，是在 $2n$ 维 Euclid 空间 $\R^{2n}$ 上取标准坐标 $(p_1,q_1,\dots,p_n,q_n)$，并由
\begin{equation}
  \omega_{\R^{2n}} = \sum_{i=1}^{n}\rd p_i\wedge\rd q_i
  \tag{2.11.27}
\end{equation}
定义辛形式而得的流形。从辛流形 $(\R^{2n},\omega_{\R^{2n}})$（或其开子流形）到自身的辛同相写像称为\textbf{正则变换}（canonical transformation）。

\takeaway{辛流形是带非退化闭 2 形式的流形，$\R^{2n}$ 为最基本之例，其自同构即正则变换。}

% [原文 p.31]
\section{探测器的轨道}

二体 Kepler 问题已由 Newton 解决，而多体问题远为困难，Euler、Lagrange、Jacobi、Poincaré 等众多数学家都曾致力于此。天体力学中的多体问题是分析力学与辛几何的故乡，同时是动力系统与混沌的源头，对数学的影响不可估量。20 世纪迎来了发射人造卫星与探测器的时代，除理论上的兴趣之外，又添了实用的动机。太空开发中事故难免，一旦发生意外，迅速找出所需时间尽可能短的轨道，或者反过来找出燃料消耗尽可能少的轨道，便成为重要课题。

\takeaway{系统显含时间则能量不守恒，探测器正因此可借引力助推加速。}

\subsection{平面圆型限制性三体问题}

三体问题研究 3 个相互施加引力的质点的运动。假定最后 1 个质点的质量比前 2 个极小，把前 2 个当作背景引力场，只研究最后 1 个质点的运动，这种近似称为\textbf{限制三体问题}（restricted 3-body problem）。太阳与木星引力场中小行星、彗星的轨道，地球与月球引力场中火箭与人造卫星的轨道，都属此类。进一步，当前 2 体的轨道为真圆时，称\textbf{圆型限制三体问题}（circular restricted 3-body problem）；再进一步，当最后 1 个质点在最初 2 体的轨道面内运动时，称\textbf{平面圆型限制三体问题}（planar circular restricted 3-body problem）。

把 2 体的质量归一化为 $\mu$ 与 $1-\mu$，并假定 $\mu\in(0,1/2]$。称较重天体为太阳、较轻天体为木星，它们的轨道分别记为 $S$ 与 $J$：
\begin{align}
  S&\colon\R\to\R^2,\quad t\mapsto(-\mu\cos t,\,-\mu\sin t),
  \tag{2.12.1}\\
  J&\colon\R\to\R^2,\quad t\mapsto\bigl((1-\mu)\cos t,\,(1-\mu)\sin t\bigr).
  \tag{2.12.2}
\end{align}

% [原文 p.32]
把太阳与木星轨道的补集
\begin{equation}
  W := \R^2\setminus\bigl(\Image(S)\cup\Image(J)\bigr)
  \tag{2.12.3}
\end{equation}
取作探测器的配位空间，则哈密顿函数
\begin{equation}
  H\colon T^*W\times\R\to\R
  \tag{2.12.4}
\end{equation}
由
\begin{equation}
  H(\bm{P},\bm{Q},t) = \frac12|\bm{P}|^2 - \frac{\mu}{|\bm{Q}-J(t)|} - \frac{1-\mu}{|\bm{Q}-S(t)|}
  \tag{2.12.5}
\end{equation}
给出。由于哈密顿函数显含时间，系统不具有时间平移不变性，能量不守恒。探测器之所以能借助引力助推（gravity assist，原书作スイングバイ）在不消耗燃料的情况下加速，原因正在于此。当然，若把太阳与木星的能量也一并计入，整个系统是能量守恒的。

\subsection{含时的正则变换}

光滑映射 $\psi\colon T^*W\times\R\to T^*W$ 满足以下两个条件时，称为\textbf{含时的正则变换}（time-dependent canonical transformation）：
\begin{itemize}[leftmargin=2em]
  \item 对任意 $t\in\R$，$\psi(-,t)\colon T^*W\to T^*W$ 是微分同胚；
  \item 存在函数 $K\colon T^*W\times\R\to\R$，使
\begin{equation}
  \psi^*\eta = \eta - \rd K
  \tag{2.12.6}
\end{equation}
成立。
\end{itemize}
其中
\begin{equation}
  \eta = \bm{p}\cdot\rd\bm{q} - H\rd t
  \tag{2.12.7}
\end{equation}
是 Poincaré--Cartan 积分不变量。作用泛函用 Poincaré--Cartan 积分不变量表为
\begin{equation}
  S\colon\gamma\mapsto\int_{t_0}^{t_1}\gamma^*\eta,
  \tag{2.12.8}
\end{equation}
因此在正则变换前后，它只发生表面项的差异：
\begin{align}
  \int_{t_0}^{t_1}\gamma^*\psi^*\eta - \int_{t_0}^{t_1}\gamma^*\eta
  &= -\int_{t_0}^{t_1}\gamma^*\rd K
  \tag{2.12.9}\\
  &= -\int_{t_0}^{t_1}\rd(\gamma^*K)
  \tag{2.12.10}\\
  &= K(\gamma(t_0)) - K(\gamma(t_1)).
  \tag{2.12.11}
\end{align}
特别地，固定端点作变分所得的哈密顿正则运动方程，在含时正则变换的前后保持不变：
\begin{equation}
  \bm{P} = \psi^*\bm{p},\qquad
  \bm{Q} = \psi^*\bm{q},\qquad
  H' = H + \frac{\partial K}{\partial t}.
  \tag{2.12.12}
\end{equation}

% [原文 p.33]
令 (2.12.12)，则 (2.12.6) 可表为
\begin{equation}
  \rd K = \bm{p}\cdot\rd\bm{q} - \bm{P}\cdot\rd\bm{Q} + (H'-H)\rd t.
  \tag{2.12.13}
\end{equation}

\subsection{Jacobi 积分与 Lagrange 点}

用旋转矩阵
\begin{equation}
  R(t) = \begin{pmatrix}\cos t & \sin t\\ -\sin t & \cos t\end{pmatrix}
  \tag{2.12.14}
\end{equation}
定义含时的坐标变换
\begin{equation}
  \bm{Q} = R(-t)\bm{q},\qquad
  \bm{Q} = \begin{pmatrix}Q_1<本地路径>,\qquad
  \bm{q} = \begin{pmatrix}q_1<本地路径>
  \tag{2.12.15}
\end{equation}
（即把 $W\times\R$ 的坐标由 $(\bm{Q},t)$ 换成 $(\bm{q},t)$）。这是线性变换，故在切空间 $T_{\bm{Q}}W$ 上引出同一映射。它在余切空间 $T^*_{\bm{Q}}W$ 上引出的映射由
\[
  \bm{p} = \bm{P}R(-t),\qquad
  \bm{P} = (P_1\ \ P_2),\qquad
  \bm{p} = (p_1\ \ p_2)
\]
给出。

由
\begin{align}
  \frac{\partial R}{\partial t}
  &= \begin{pmatrix}-\sin t & \cos t\\ -\cos t & -\sin t\end{pmatrix}
  \tag{2.12.16}\\
  &= \begin{pmatrix}\cos t & \sin t\\ -\sin t & \cos t\end{pmatrix}
     \begin{pmatrix}0 & 1\\ -1 & 0\end{pmatrix}
  \tag{2.12.17}\\
  &= R(t)\begin{pmatrix}0 & 1\\ -1 & 0\end{pmatrix}
  \tag{2.12.18}
\end{align}
可得
\begin{align}
  \rd\bm{Q} &= \rd\bigl(R(-t)\bm{q}\bigr)
  \tag{2.12.19}\\
  &= R(-t)\rd\bm{q} - \frac{\partial R}{\partial t}(-t)\,\bm{q}\,\rd t
  \tag{2.12.20}\\
  &= R(-t)\left(\rd\bm{q} - \begin{pmatrix}0 & 1\\ -1 & 0\end{pmatrix}\bm{q}\,\rd t\right),
  \tag{2.12.21}
\end{align}
因此经此变换，Liouville 形式变为
\begin{align}
  \bm{P}\cdot\rd\bm{Q}
  &= \bm{P}R(t)\left(\rd\bm{q} - \begin{pmatrix}0 & 1\\ -1 & 0\end{pmatrix}\bm{q}\,\rd t\right)
  \tag{2.12.22}\\
  &= \bm{p}\cdot\rd\bm{q} - \mu\,\rd t.
  \tag{2.12.23}
\end{align}

% [Tutor Note: 校勘——(2.12.22) 原书作 $\bm{P}R(t)$，但由 (2.12.21) 与 $\bm{p}=\bm{P}R(-t)$，此处应为 $\bm{P}R(-t)$；译文照录原书。]

% [原文 p.34]
其中
\begin{align}
  \mu &= \bm{p}\begin{pmatrix}0 & 1\\ -1 & 0\end{pmatrix}\bm{q}
  \tag{2.12.24}\\
  &= (p_1\ \ p_2)\begin{pmatrix}0 & 1\\ -1 & 0\end{pmatrix}\begin{pmatrix}q_1<本地路径>
  \tag{2.12.25}\\
  &= p_1q_2 - p_2q_1
  \tag{2.12.26}
\end{align}
是生成旋转的哈密顿函数。由 (2.12.13)，Poincaré--Cartan 积分不变量变为
\begin{equation}
  \bm{P}\cdot\rd\bm{Q} - H\rd t = \bm{p}\cdot\rd\bm{q} - \mu\rd t - H\rd t = \bm{p}\cdot\rd\bm{q} - h\rd t,
  \tag{2.12.27}
\end{equation}
故旋转坐标系 $(\bm{p},\bm{q})$ 中的哈密顿函数由
\begin{equation}
  h = H + \mu = \frac12\bigl(p_1^2+p_2^2\bigr) - \frac{\mu}{|\bm{q}-J(0)|} - \frac{1-\mu}{|\bm{q}-S(0)|} + p_1q_2 - p_2q_1
  \tag{2.12.28}
\end{equation}
给出。注意变换后的哈密顿函数 $h$ 不依赖时间。于是 $h$ 成为系统的守恒量，称为 \textbf{Jacobi 积分}（Jacobi integral）。对 (2.12.28) 配方可得
\begin{equation}
  h = \frac12\Bigl((p_1+q_2)^2+(p_2-q_1)^2\Bigr) - \frac{\mu}{|\bm{q}-J(0)|} - \frac{1-\mu}{|\bm{q}-S(0)|} - \frac12|\bm{q}|^2.
  \tag{2.12.29}
\end{equation}
哈密顿函数中不依赖 $\bm{p}$ 的部分
\begin{equation}
  U\colon\R^2\setminus\{J(0),S(0)\}\to\R,\qquad
  \bm{q}\mapsto -\frac{\mu}{|\bm{q}-J(0)|} - \frac{1-\mu}{|\bm{q}-S(0)|} - \frac12|\bm{q}|^2
  \tag{2.12.30}
\end{equation}
称为\textbf{有效势}（effective potential）。

本系统的平衡点是哈密顿函数 $h$ 的微分为零的点，称为 \textbf{Lagrange 点}（Lagrangian point）。关于 $\bm{p}$ 的微分为零的条件给出
\begin{equation}
  p_1 = -q_2,\qquad p_2 = q_1.
  \tag{2.12.31}
\end{equation}
在此条件下，关于 $\bm{q}$ 的微分为零的条件对应于有效势的临界点。直接计算可知，Lagrange 点共有 5 个：前 3 个位于连接太阳与木星的直线上，其余 2 个（\textbf{Trojan 点}（Trojan points））分别与太阳、木星构成正三角形。对给定的实数 $c$，能量的等位集合（level set）$h^{-1}(c)$ 称为\textbf{能量超曲面}（energy hypersurface），其在 $W$ 上的射影之像称为 \textbf{Hill 区域}（Hill's region）。

% [原文 p.35]
Hill 区域是探测器在给定能量下所能到达的太阳系内的区域。由哈密顿函数 (2.12.29) 的第 1 项恒正这一事实可知，Hill 区域由
\begin{equation}
  \{\bm{q}\in\R^2\setminus\{S(0),J(0)\}\mid U(\bm{q})\le c\}
  \tag{2.12.32}
\end{equation}
给出。

% [Tutor Note: 背景——Hill 区域得名于美国天文学家 G. W. Hill（1838--1914）的月球运动研究，其边界即零速度曲线。]

\chapterreview{本章由 Newton 力学出发，先建立 Lagrange 形式，再过渡到相空间上的 Hamilton 形式，并把正则方程改写为泊松括号的形式。泊松括号的反对称性、Leibniz 法则与 Jacobi 恒等式被抽象为泊松流形的定义，哈密顿向量场随之成为由 $\{f,-\}$ 确定的向量场。Noether 定理指出，哈密顿流的对称性与系统的守恒量是同一件事，能量守恒与动量守恒都只是它的特例。辛流形是泊松流形在非退化条件下的特殊情形，其辛叶即泊松结构的叶，辛同相写像则保持辛形式不变。最后以平面圆型限制性三体问题为例，说明在显含时间的哈密顿系统中，Jacobi 积分、Lagrange 点、Hill 区域与引力助推如何共同刻画探测器的轨道。}
